\section{Cryptography Basics}\label{sec:basics}
  \todoenum[rough structure]{
    \item PKE and KEM: Definitions, OW vs IND, CPA vs CCA
    \item Signatures (and NIZKs?): Definition, EU-CMA
    \item Game hopping
  }

  \subsection{Sets, Lists, and Probabilities}
    Let us start by defining our conventions. 
      We will take the natural numbers $\NN$ to include $0$, and for $n \in \NN$, we
      will use $[n]$ to denote the set $\{0, \dots, n-1\}$. We will also use
      capital letters to denote powers of two, such as $N = 2^n$ (and $NM =
      2^{n+m}$).

    Functions are defined over bitstrings unless otherwise noted, and the set
      of all functions from the set of $m$-length bitsrings $\bin^m$ to the
      set of $n$-length $\bin^n$ bitstrings is denoted $\funcspace{m}{n}$.
      The $n$-by-$n$ identity matrix is denoted $\idmatrix[n]$, where we'll
      drop the subscript whenever dimension is clear from context.

    For a bitstring $x\in\{0,1\}^*$, we let $\card{x}$ denote its length;
      for a finite set $\Xspace$, we let $\card{\Xspace}$ denote its
      cardinality; and likewise for lists.
      For a list $\Xlist$, $\Xlist[i]$ retrieves its $i$th element,
      and two lists $\Xlist$ and $\Ylist$ may be combined by concatination to form
      a new list $\Zlist = \Xlist \concat \Ylist$. It will at times be useful
      to treat bitstrings as lists of binary values.

    Deterministic assignment to a (classical) variable is denoted $x \gets y$;
      uniform sampling from a set $\Xspace$ (and assignment of the result to
      $x$) is denoted $x \sample \Xspace$. Likewise, assignment of outputs from
      deterministic resp.\ probabilistic algorithms $X$, $Y$
      is denoted $x \gets X(\cdot)$ and $y \sample Y(\cdot)$, respectively.

    We use $\condprobsample{\mathsf{Code}}{\mathsf{Event}}{\mathsf{Condition}}$
      to denote the conditional probability of $\mathsf{Event}$ occurring
      when $\mathsf{Code}$ is executed, conditioned on $\mathsf{Condition}$
      (where the probability is taken over the random coins of
      $\mathsf{Code}$).
      We omit $\mathsf{Code}$ when it is clear from context and
      $\mathsf{Condition}$ when it is not needed. 

  \subsection{Algorithms and Oracles}
    Algorithms are assumed to be either classical~\cite{Turing36} or quantum Turing
      machines~\cite{Deutsch85} (although we will follow the convention of
      expressing quantum algorithms as quantum circuits~\cite{Yao93});
      which of the two we are talking about will usually be clear from context.
      An \emph{oracle} algorithm is an algorithm $A$ that additionally has black-box access to
      some oracle $\Oracle$, meaning it can only call the oracle on some input
      and receive its output; this is written $A^\Oracle$. 

    Unlike algorithms, oracles are typically not limited to bounded resources.
      Classical oracles take a classical value $x$ as input and return some
      value $y$, and can be deterministic (repeat queries return the same value) or
      probabilistic (the output depends on randomness internal to the oracle);
      in either case we write $y \gets \Oracle(x)$.
      Oracles can furthermore be either stateful (queries are remembered) or
      stateless. Clearly, stateless, deterministic oracles implement a function,
      and we may denote an oracle implementing the function $f$ as $\Oracle_f$.

    Quantum(-accessible) oracles\footnote{
      Some works distinguish between quantum-\emph{accessible} oracles,
      which evaluate a function on a superposition of inputs, and
      \emph{quantum} oracles, which apply some quantum operation on the input
      registers.
      We do not distinguish between these two kinds of oracle, and rather view
      the former as a special case of the latter.
      } are oracles that instead take a set of \emph{qubit registers} as input,
      apply some operation on them, before returning the registers to the
      querying party. The operation applied may be probabilistic in nature and
      may include measurements (which has the effect of transforming pure
      quantum states into mixed states\footnote{
        Such quantum operators are known as \emph{superoperators}
      }), though it will often be required to be unitary.
      For quantum oracle algorithms, we will take a classical oracle to be a
      quantum oracle that performs a computational-basis measurement on the
      registers before operating on them. Hence, we do not make a
      distinction between classical and quantum oracle access 
      in our notation, writing $A^\Oracle$ for either case. 

    In \autoref{subsec:pseudocode}, we define a quantum pseudocode that allows
      us to talk about general quantum oracles, which we apply to quantum random
      oracles in \autoref{sec:rom}. For now, we note that a quantum oracle
      $\Oracle_f$ that evaluates a function $f$ on superpositions of inputs are
      defined as the oracle that applies the unitary operation
      \[
        \ket{x,y} \to \ket{x,y \xor f(x)}\, ,
      \]
      where $\ket{x,y}$ are the computational basis states of the qubit
      input/output registers given to the oracle. (See \autoref{sec:primer} for
      an introduction to the basics of quantum computation.)

    %% MORE FROM SOA SOK
    %We use code-based experiments, where by convention all sets, lists,
      %and lazily sampled functions are initialized empty.
      %In our experiments, we will assume that no oracle calls are made (by the adversary)
      %with evidently out-of-bounds inputs (e.g.~list indices or key handles).

    %We use the shorthand $\Xset \setapp x$ for
      %the operation $\Xset \gets \Xset \cup \{ x \}$ and
      %$\Xlist \listapp x$ for appending the element $x$ to a list $\Xlist$.

    %We can make the
      %randomness $\rand$ of $Y$ explicit by writing
      %$x \gets Y(\, \cdot\, ; \rand)$, for previously sampled $\rand$.

    %We will also consider stateful algorithms, for instance when we write
    %  $\mssg \sample \mssgsamp[\samplestate](\dist)$
    %  the algorithm $\mssgsamp$ takes as state $\samplestate$ and
    %  as input the value $\dist$ and outputs $\mssg$,
    %  but it may change its state $\samplestate$ as part of its processing
    %  (code-snippet taken from \autoref{fig:IND-SOA}).

  \subsection{Public Key Encryption}
    First, whenever some object is generated
      dependent on some other public object, such as the public key $\pk$ being
      dependent on the public parameters $\params$ (which in turn depend on the
      security parameter $\secpar$), we assume without loss of generality that
      the new object includes descriptions of the prior (public) objects (so if you
      know $\pk$, then you also know $\params$). Similarly, we assume
      that if you know a secret key $\sk$ belonging to a keypair $(\pk, \sk)$, then 
      you also know $\pk$. This is w.l.o.g.\ in the sense
      that after a keypair $(\pk,\sk)$ is generated, one can always define the
      actual keypair to be $(\pk, (\pk, \sk))$; but it is not w.l.o.g.\ in
      the sense that this enfores that secret keys have unique public keys.\footnote{
        There exist PKE schemes that violate this, and
        there are even security proofs that rely on such non-unique key
        relations~\cite{EC:BelHofYil09}.
      }

    Adversaries are assumed to be generally quantum agents, having
      access to as much classical and quantum memory as they need, on which
      they can apply unitary operations and measurements as they
      please. (We will come upon cases later where it will be beneficial to restrict this
      assumption.) We will mostly think of these agents as being \emph{efficient} in the sense of
      running in time polynomial in the security parameter $\secpar$, although
      we will stick with a concrete formal treatment wherein asymptotics are
      avoided, and the running time of reductions are instead stated directly in terms
      of the running time of the original adversary. We may also
      consider \emph{unbounded} adversaries; these can be thought of as running
      in time (and space) exponential in the security parameter.

    If the adversary $\adv$ has access to a random oracle $\RO$, this is
      written $\adv^{\RO}$, and likewise for any other oracles it may have
      access to. (Our notation does not distinguish between quantum and
      classical oracles.) Similarly, if a reduction $\bdv$ has black-box access
      to an adversary $\adv$ (meaning it may control its input and read its output, but not
      otherwise read or affect its internal workings), this is written $\bdv^\adv$.

    \subsubsection{Defining security.}
      The standard security notion for public-key encryption is usually taken to
        be ``Indistinguishability under Chosen-Ciphertext Attacks'', or $\indcca$,
        but for the sake of this tutorial we will stick to the weaker 
        ``Indistinguishability under Chosen-\emph{Plaintext} Attacks'' ($\indcpa$). Since
        this is a public-key encryption scheme, the chosen-plaintext attacks can
        be performed ``offline'' by the adversary simply by running the
        encryption algorithm itself using the provided public key. 
        %Therefore, unlike in the symmetric encryption case, we do not need to provide 
        %the adversary with an honest encryption via an oracle.
        % Eh, it gets this anyway, by letting m0 = m1, making the point moot.

      On the left of \autoref{fig:ind-ow} is the experiment
        defining our (single-user) $\indcpa$ security notion. The game starts by sampling
        parameters $\params$, a keypair
        $(\pk,\sk)$, and a challenge bit $\chalbit$, before handing the public
        key to the adversary $\adv$. If the construction in question uses a
        random oracle $\RO$, then $\adv$ also gets oracle access to $\RO$. 

      At some point, the adversary decides on two equal-length messages
        $\mssg_0$ and $\mssg_1$, and call its \emph{challenge encryption} oracle $\Oenc$
        (\autoref{fig:ind-ow}, middle). This oracle then provides an honest
        encryption of the message $\mssg_{\chalbit}$ and returns a challenge
        ciphertext $\chalctxt$. If the adversary ever calls
        the challenge oracle on messages of differing lengths, then it returns the
        special ``fail'' symbol $\fail$ instead; this is due to the
        assertion that a scheme should be considered
        secure even if ciphertexts leak information about message lengths.
        In fact, you may notice from \autoref{fig:bre} that
        our BRE scheme leaks the message length exactly, since $\demctxt$
        must necessarily have the same length as $\mssg$. 
        
      We are in the \emph{multi-challenge} setting, and so the adversary may
        continue to make as many calls to $\Oenc$ as it likes,
        in addition to arbitrary calls to the random oracle $\RO$. 
        (If the adversary is restricted to a single call to $\Oenc$, we call this
        the \emph{single-challenge setting}.)
        Once it is satisfied, $\adv$ halts with a guess $\advbit \in
        \bin$, and wins if it guessed $\chalbit$ correctly.

      We measure security in terms of $\adv$'s \emph{distinguishing advantage},
        i.e.\ how much better it can do than random guessing.

      \begin{figure}[t]
        \centering
        \begin{pchstack}[center]
          \procedure{$\experiment{\indcpa}{\pkescheme}$}{
            \params \sample \pkpm(\secparam) \\
            (\pk, \sk) \sample \pkkg(\params) \\
            \chalbit \sample \bin \\
            \advbit \sample \adv^{\Oenc(,\RO)}(\pk) \\
            \pcreturn \advbit = \chalbit
          }
    
          \pchspace
    
          \procedure{$\indOenc(\mssg_0, \mssg_1) \pccomment{\ind}$}{
            % Single-challenge:
              %\pcif \chalctxt \neq \emptystring \lor \abs{\mssg_0} \neq \abs{\mssg_1} \pcthen \\ 
            \pcif \abs{\mssg_0} \neq \abs{\mssg_1} \pcthen \pcreturn \fail \\
            \chalctxt \sample \pkenc_\pk(\mssg_\chalbit) \\
            \pcreturn \chalctxt
          }
    
          \pchspace
    
          \procedure{$\experiment{\owcpa}{\kemscheme}[(\adv)]$}{
            \params \sample \kempm(\secparam) \\
            (\pk, \sk) \sample \kemkg(\params) \\
            i \gets 0 \\
            (i,\advdemkey_i) \sample \adv^{\Oenc(,\RO)}(\pk) \\
            \pcreturn (\advdemkey_i = \chalkeyi)
          }
    
          \pchspace
    
          \procedure{$\owOenc \pccomment{\oneway}$}{
            i \gets i+1 \\
            (\chalkeyi, \chalctxt) \sample \kemenc_\pk \\
            \pcreturn \chalctxt
          }
        \end{pchstack}
      \caption{
        \label{fig:ind-ow}
        The experiments defining (multi-challenge) $\indcpa$ for PKE (left), and
        $\owcpa$ for KEM (right). If a random oracle $\RO$ is present, then the
        adversary gets access to this in addition to the challenge oracle $\Oenc$.
      }
      \end{figure}

      \begin{definition}[Distinguishing Advantage]
        The distinguishing advantage of an adversary $\adv$ against a
          public-key encryption scheme $\pkescheme$, under chosen ciphertext
          attacks, is defined as
        \begin{align*}
          \advantage{\indcpa}{\pkescheme} &\coloneqq \abs{
            2 \cdot \prob{\experiment{\indcpa}{\pkescheme}} - 1
          }\, ,
          \intertext{or, equivalently,}
          \advantage{\indcpa}{\pkescheme} &\coloneqq \abs{
            \prob{\experiment{\indcpa}{\pkescheme}} 
            - \prob{\neg\experiment{\indcpa}{\pkescheme}}
          }\, ,
          \intertext{or, equivalently,}
          \advantage{\indcpa}{\pkescheme} &\coloneqq \abs{
            \condprobsample{
              \experiment{\indcpa}{\pkescheme}
            }{0 \sample \adv}{\chalbit = 0}
            - \condprobsample{
              \experiment{\indcpa}{\pkescheme}
            }{0 \sample \adv}{\chalbit = 1}
          }\, .
        \end{align*}
        where $\experiment{\indcpa}{\pkescheme}[]$ is as given in
        \autoref{fig:ind-ow}.
      \end{definition}

      \noindent
      As is easiest seen from the first of the above three equations, the
        advantage is just the probability that $\adv$ wins the game, minus the
        probability of winning based on a random guess, scaled so that the
        advantage is a value defined on the interval from $0$ to $1$,
        where $0$ means $\adv$ has no advantage (e.g.\ if the messages are perfectly hidden), 
        and $1$ means that $\adv$ is a perfect distinguisher (the scheme is
        completely broken, or $\adv$ is unbounded).

  \subsection{Key Encapsulation Mechanisms}

  \subsection{Digital Signature Algorithms}

  \subsection{Security Reductions}
    \subsubsection{Kinds of reductions.}
      We say that a reduction is \emph{simple} if it only runs the adversary once,
        and \emph{fully black-box} if it runs the adversary from start to finish,
        answering oracle queries along the way, and otherwise only providing its
        input and reading its output; if, as is typically the case, the reduction
        is both simple and fully black-box, we call it an SFBB reduction. 

      We will later see some reductions that \emph{rewind} the adversary, and run
        it again from a point chosen by the reduction. We model this as the usual
        black-box access with an additional rewinding \emph{interface}, with which
        it can essentially tell the adversary to rewind itself to a specified point
        in its computation; we refer to it as a \emph{rewinding black-box} (RBB)
        reduction.


    \subsubsection{Game Hopping.}
      %\todoenum[Provide explanations of:]{
      %  \item Basic background on game hopping and advantages, stuff like
      %    difference vs.\ times two minus 1, with or
      %    without absolute values, the fundamental lemma of game hopping.
      %  \item Maybe write something about asymptotic vs.\ concrete security.
      %}
      For our security proofs, we will be working in the game-hopping
        paradigm of Shoup~\cite{EPRINT:Shoup04}.
        The fundamental lemma of game hopping states that the distinguishing advantage
        of an adversary can be bounded in terms of their distinguishing advantage
        over several intermediary game ``hops''. 

      \begin{lemma}[Fundamental Lemma of Game Hopping]\label{lemma:game-hopping}
        For any adversary $\adv$ tasked with distinguishing $\game{0}$ from
        $\game{n}$, and for any games $\game{1}, \ldots \game{n-1}$,
        there exist adversaries $(\adv_1, \ldots, \adv_n)$ such that
        \[ 
          \abs{
            \probsample{\game{0}(\adv)}{1 \sample \adv} - 
            \probsample{\game{n}(\adv)}{1 \sample \adv}
          } \leq \sum_{i = 1}^n \abs{
            \probsample{\game{i-1}(\adv)}{1 \sample \adv_i} - 
            \probsample{\game{i}(\adv)}{1 \sample \adv_i}
          }\, .
        \]
      \end{lemma}

      \noindent
      The lemma states that, if $\adv$ were able to distinguish, then
        at least one of $\adv_1, \ldots, \adv_n$ must be able to
        distinguish too---and vice versa, that if none of the $\adv_i$ are able to
        distinguish, then neither is $\adv$.
        (In fact, the $\adv_i$ that achieve this will in essence be just copies
        of $\adv$.)

      The proof is remarkably simple, reproduced here for those who haven't 
        seen it.

      \begin{proof}
        For all $i$ from $1$ to $n-1$, insert $0 = \probsample{\game{i}(\adv_i)}{1 \sample \adv_i} -
        \probsample{\game{i}(\adv_i)}{1 \sample \adv_i}$ into the absolute value
        on the left hand side, then apply the union bound.
        \qed
      \end{proof}

    %\paragraph{Quantum notation.}
    %  %  While we only assume knowledge of basic linear algebra, this tutorial 
    %  %    still concerns itself with several fundamentally quantum concepts. To
    %  %    help bridge this divide, we provide a self-contained primer on quantum
    %  %    computing theory in
    %  %    \autoref{sec:primer}, with explanations and definitions for all
    %  %    the quantum concepts and conventions that we will need. A summary of our
    %  %    notation is given in \autoref{tab:notation}.
    %  Precise definitions for our quantum-notational conventions are provided in
    %    \autoref{sec:primer},
    %    along with a gentle introduction to the relevant quantum concepts,
    %    summarized in \autoref{tab:notation}.

      \hhnote{
        Add: $\pcparse$ keyword. (Allow $\pcif \pcparse$ too? By having $\pcparse$
        return a boolean value? Hmm, tricky.) And: $\red\bad$ abort convention.
      }
      \hhnote{
        Add: Adversaries are always assumed to halt in due time, even when the
        environment fails to match their expectation.
      }

    \subsubsection{Bad Events.}
      \hhnote{
        Add identical-until-bad lemma?
      }
